% Full-parameter training: random seed, conditional process theorem, exact matrix dynamics,
% conditional cancellation, open obligations and diagnostics. Input from appendix.tex.

\section{Training every parameter}\label{app:full}
This appendix proves the statements of \cref{sec:full}. It does not modify \cref{thm:step} or \cref{cor:matched}. Throughout, every tensor of
\cref{eq:block} (token and position embeddings, the $18$ query, key and value maps, the two MLPs with their biases, the readout and
the relative biases) follows $\dot\Theta=-\nabla_\Theta\mathcal L$ in the Euclidean/Frobenius metric.

\subsection{The random seed}\label{app:seed}
\begin{proof}[Proof of \cref{lem:seed}]
Write $\Theta_0=\varepsilon\tilde\Theta$ with $\tilde\Theta=G+(\delta/\varepsilon)A$. Let $\mathrm B$ be a Gaussian block of size $d_1\times d_2$ and
$m=\max(d_1,d_2)$, so $\sqrt m\,G_{\mathrm B}$ has i.i.d.\ standard normal entries. By the Gaussian singular-value bounds
\citep[Chapter~7]{vershynin2018high}, for every $t\ge0$, with probability at least $1-2e^{-t^2/2}$,
\[
 \sqrt{m}-\sqrt{\min(d_1,d_2)}-t\ \le\ \sigma_{\min}(\sqrt m\,G_{\mathrm B})\ \le\ \sigma_{\max}(\sqrt m\,G_{\mathrm B})\ \le\ \sqrt{d_1}+\sqrt{d_2}+t .
\]
\emph{Operator norms.} Take $t=\sqrt m$: then $\|G_{\mathrm B}\|_{\mathrm{op}}\le3$ with probability at least $1-2e^{-m/2}$.
\emph{Embeddings.} The token embedding has $d_{\mathrm{res}}\ge16|\mathcal V|$ rows. Take $t=\sqrt{d_{\mathrm{res}}}/4$: then
$\sigma_{\min}(G_E)\ge1-\frac14-\frac14=\frac12$ with probability at least $1-2e^{-d_{\mathrm{res}}/32}$. Restricting to the columns of the value
tokens cannot decrease the smallest singular value of a tall matrix.
\emph{Union bound.} Since $m\ge m_{\min}$ for every block and $d_{\mathrm{res}}\ge m_{\min}$, all these events hold together with probability
at least $1-2Ne^{-m_{\min}/2}-2e^{-m_{\min}/32}\ge1-2(N+1)e^{-m_{\min}/32}$.
\emph{Perturbation.} $\|(\delta/\varepsilon)A_{\mathrm B}\|_{\mathrm{op}}\le\frac14$, so by Weyl's inequality every scaled block has operator norm at most
$\frac{13}{4}\le a_0=4$ and the value embedding has $\sigma_{\min}\ge\frac12-\frac14=c_{\mathrm{emb}}$.
\emph{Biases, readout, relative biases.} The scaled MLP biases are signs, so each coordinate has magnitude $1$ and each vector has norm
$\sqrt w=N_b$; $A$ vanishes on them. The readout and the relative biases are zero. Hence \cref{ass:init} holds with
$r_0=\min(\frac18,\frac{c_{\mathrm{emb}}}4)=\frac1{16}$.
\end{proof}

\begin{proof}[Proof of \cref{prop:random-plateau}]
On the event of \cref{lem:seed} the seed satisfies \cref{ass:init}, has zero readout and $\beta=0$, and the masks are full causal masks.
The ambient flow argument at the end of Appendix~\ref{app:outcome} applies with $P(t)=I$, i.e.\ every parameter trained: the loss floor
$\Lout\ge\log K-C_2\varepsilon^2\lambda^n$ holds on the enlarged ball, and the exit time from radius $r_0\varepsilon$ is at least
$r_0^2\lambda^{-n}/C_2$. On $[0,T_{\mathrm{out}}^{\mathrm{rand}}]$, half of that, the displacement is at most $r_0\varepsilon/\sqrt2$. The upper bound
$\Lout\le\log K$ holds because the loss starts at $\log K$ and does not increase. For the probability,
$|\E p_{t,y}-\frac1K|\le\E|p_{t,y}-u_y|\le\E\,\mathrm{TV}(p_t,u)\le\sqrt{\E\KL(u\|p_t)/2}\le\sqrt{C_2/2}\,\varepsilon\lambda^{n/2}$ by Pinsker and
Jensen. No metric comparison is needed, since the flow is already Euclidean on all parameters.
\end{proof}

\paragraph{Why the readout starts at zero and the biases are signs.} With a random readout of scale $\varepsilon$, $\Lout(\Theta_0)-\log K$
is of order $\varepsilon^2$, independent of $n$; the displacement bound then gives only $\sqrt{t\varepsilon^2}$ and the exponential interval is lost.
Gaussian biases would have coordinates arbitrarily close to zero, so no neighbourhood of the seed would have a fixed ReLU pattern. Both
choices are part of the deterministic seed, not of $G$.

\subsection{Proofs of \texorpdfstring{\cref{thm:process,prop:separable}}{Theorem 2 and Proposition 3}}\label{app:process}
\begin{proof}[Proof of \cref{thm:process}]
Write $L(t)=\Lproc(\Theta(t))$ and $c=a_n^2/B_n^2$.
\emph{Gradient domination.} While $L>\tau$, Cauchy--Schwarz and \ref{hyp:align} give
$a_nL\le-\langle\nabla L,D\rangle\le\|\nabla L\|B_n$, hence $\|\nabla L\|^2\ge cL^2$.
\emph{Decay.} By \ref{hyp:regular}, $\dot L=-\|\nabla L\|^2\le-cL^2$ almost everywhere while $L>\tau$, so $\frac{d}{dt}(1/L)\ge c$ and
\[
 L(t)\le\frac{L_0}{1+cL_0t}<\frac1{ct}.
\]
Let $t^\ast$ be the first time with $L\le\tau$; then $t^\ast\le1/(c\tau)=T_{\mathrm{proc}}$. Since $L$ is nonincreasing, $L(t)\le\tau$ for all $t\ge t^\ast$.
\emph{Teacher-forced errors.} Let $E_m$ be the event that greedy decoding on the gold prefix $(P,y_{<m})$ does not return $y_m$. If greedy
returns $k\ne y_m$, then $p_\Theta(y_m\mid\cdot)\le p_\Theta(k\mid\cdot)$, so $p_\Theta(y_m\mid\cdot)\le\frac12$ and the token loss is at least $\log2$.
By Markov's inequality, $\Pr_P(E_m)\le\E[-\log p_\Theta(y_m\mid P,y_{<m})]/\log2$.
\emph{First error.} If free-running decoding errs, let $m$ be its first wrong token. All earlier generated tokens are correct, so the model
sees exactly the gold prefix at $m$ and $E_m$ occurs. Therefore
\[
 \Pr_P[\text{the continuation is wrong}]\le\sum_{m=1}^M\Pr_P(E_m)\le\frac{M\,L(t)}{\log2}\le\frac{M\tau}{\log2}=\xi .
\]
Under \ref{hyp:poly} and $\xi^{-1}\le\mathrm{poly}(n)$, $T_{\mathrm{proc}}=MB_n^2/(a_n^2\xi\log2)$ is polynomial, since $M=n+3$.
\end{proof}

\begin{proof}[Proof of \cref{prop:separable}]
For one context, $\nabla_z\ell=p-e_y$ and $z=Vh$, so $\nabla_V\ell=(p-e_y)h^\top$ and
\[
 -\langle\nabla_V\ell,U\rangle=-\sum_kp_kU_k^\top h+U_y^\top h=\sum_{k\ne y}p_k\,(U_y-U_k)^\top h\ \ge\ \gamma_n(1-p_y)
\]
by \ref{hyp:sep}. With $\ell=-\log p_y$, $1-p_y=1-e^{-\ell}\ge\ell/(1+\ell)\ge\ell/(1+H_n)$ by \ref{hyp:lossbd}. Averaging over contexts and positions
gives $-\langle\nabla_\Theta\Lproc,D\rangle=-\langle\nabla_V\Lproc,U\rangle\ge\frac{\gamma_n}{1+H_n}\Lproc$, and $\|D\|=\|U\|_F\le B_n$.
Substituting $a_n$ into \cref{eq:proc-time} gives \cref{eq:sep-time}. The witness is not a training constraint: $\Phi$ still follows its
own gradient.
\end{proof}

\paragraph{How strong is \ref{hyp:sep} at the seed?} With the readout at zero, $z=Vh$ has no derivative in $\Phi$, so $\nabla_\Phi\Lproc(\Theta_0)=0$
and the first motion is readout-only. Near the seed the ReLU pattern is fixed, the network at zero attention scores is affine in the
individual tokens, and the score-dependent terms are smaller by a factor $O(\varepsilon^4)$ (\cref{prop:score-degree}). A readout of features that
are additive in $(s_{t-1},b_t)$ cannot give the label $s_{t-1}+b_t$ a positive margin: with $c_k$ the change of logit $k$ when $b_t$ flips
from $0$ to $1$, each $s$ requires $c_{s+1}-c_s>2\gamma_n$, and these differences sum to zero around the cycle. If $\delta$ bounds $|(U_a-U_b)^\top(h-h^{[0]})|$, with $h^{[0]}$ the
zero-score features, the additive part has margins at least $\gamma_n-\delta$, so the cyclic sum forces $\gamma_n\le\delta\le\sqrt2\,\|U\|_F\sup\|h-h^{[0]}\|$:
at the seed \ref{hyp:sep} holds at best with a margin of score order. The readout witness is therefore mainly a certificate for the phase after an interaction feature has formed. The early phase
needs a witness $D$ with components outside the readout, or a separate bootstrap argument.

\subsection{Exact full-matrix dynamics}\label{app:matrix}
Fix one head with query and key maps $W_Q,W_K\in\R^{d_{\mathrm h}\times d_{\mathrm{res}}}$. Its scores are
$h_i^\top W_Q^\top W_Kh_j/\sqrt{d_{\mathrm h}}$, so the loss depends on $(W_Q,W_K)$ only through $M=W_Q^\top W_K$. Put $N=\nabla_ML$.

\begin{proposition}[Query--key flow]\label{prop:qk-flow}
Along the gradient flow of any loss, with all other tensors also moving,
\begin{equation}\label{eq:qk-flow}
 \begin{gathered}
 \dot W_Q=-W_KN^\top,\qquad \dot W_K=-W_QN,\qquad \dot M=-\mathcal P(N),\\
 \mathcal P(Z):=(W_Q^\top W_Q)Z+Z(W_K^\top W_K),
 \end{gathered}
\end{equation}
and $W_QW_Q^\top-W_KW_K^\top$ is constant in time. $\mathcal P$ is self-adjoint and positive semidefinite in the Frobenius inner product, with
$\langle Z,\mathcal PZ\rangle=\|W_QZ\|_F^2+\|ZW_K^\top\|_F^2$.
\end{proposition}
\begin{proof}
$dL=\langle N,dW_Q^\top W_K+W_Q^\top dW_K\rangle=\langle W_KN^\top,dW_Q\rangle+\langle W_QN,dW_K\rangle$ gives the first two equations, and
$\dot M=\dot W_Q^\top W_K+W_Q^\top\dot W_K$ the third. Then $\frac{d}{dt}W_QW_Q^\top=-W_KN^\top W_Q^\top-W_QNW_K^\top=\frac{d}{dt}W_KW_K^\top$. The last claims follow
from $\langle Y,\mathcal PZ\rangle=\operatorname{tr}(Y^\top W_Q^\top W_QZ)+\operatorname{tr}(Y^\top ZW_K^\top W_K)$.
\end{proof}

\begin{proposition}[Score sensitivity through the ReLU MLP]\label{prop:score-sens}
In the second block, let $\bar h_i$ be the post-attention residual at position $i$, $v_j$ the value written from site $j$ by one head,
$\omega_{ij}$ its weights, $\bar v_i=\sum_k\omega_{ik}v_k$, $D_i=\operatorname{diag}\mathbf 1[W_{\mathrm{in}}\bar h_i+b>0]$, and $r\in\R^{d_{\mathrm{res}}}$ any readout
functional, such as a row difference $V_y-V_k$. Off the ReLU switching set,
\begin{equation}\label{eq:score-sens}
 \frac{\partial\,r^\top h_i^2}{\partial s_{ij}}=\omega_{ij}\,\bigl[(I+W_{\mathrm{in}}^\top D_iW_{\mathrm{out}}^\top)r\bigr]^\top(v_j-\bar v_i).
\end{equation}
\end{proposition}
\begin{proof}
$\partial\omega_{ik}/\partial s_{ij}=\omega_{ij}(\mathbf 1[k=j]-\omega_{ik})$, so $\partial\bar h_i/\partial s_{ij}=\omega_{ij}(v_j-\bar v_i)$, and
$\partial h_i^2/\partial\bar h_i=I+W_{\mathrm{out}}D_iW_{\mathrm{in}}$.
\end{proof}
For a first-block head the same formula holds with $r$ replaced by its pullback through the Jacobian of the second block. The three identities and the conservation law were also checked against automatic differentiation on random instances (\texttt{verification/full\_matrix\_checks.py}).

\paragraph{Signal and interference.} Fix a reference matrix $S$ in a chosen task subspace $\mathcal T$ (for example, the score directions that
move mass to the displayed state or the current letter) and split the descent direction as $-N=\Pi_{\mathcal T}(-N)+E$. By
\cref{eq:qk-flow},
\[
 \frac{d}{dt}\langle S,M\rangle=\langle S,\mathcal P\,\Pi_{\mathcal T}(-N)\rangle+\langle S,\mathcal PE\rangle .
\]
If $S$ moves in time, add $\langle\dot S,M\rangle$. Taking $S=\Pi_{\mathcal T}(-N)$ at a given instant, a sufficient condition for the task score to
grow at half its interference-free rate is
\begin{equation}\label{eq:interference}
 \langle S,\mathcal PE\rangle\ \ge\ -\tfrac12\langle S,\mathcal PS\rangle .
\end{equation}
The restricted family of \cref{thm:step} has no interference by construction; full matrices make $E$ explicit. The conserved $W_QW_Q^\top-W_KW_K^\top$ is the matrix form of the scalar balance $q^2-k^2$ of \cref{rem:J-balance}.

\subsection{A conditional Gaussian cancellation lemma}\label{app:cancel}
\begin{lemma}[Conditional cancellation]\label{lem:cancel}
Let $\mathcal F$ be a $\sigma$-field, $S$ and $\mathcal P$ be $\mathcal F$-measurable with $\mathcal P$ self-adjoint and positive semidefinite, and
$E=\sum_ig_iE_i$ with $E_i$ $\mathcal F$-measurable and $(g_i)$ i.i.d.\ standard normal, independent of $\mathcal F$. Let
$\Sigma_E=\sum_i\operatorname{vec}(E_i)\operatorname{vec}(E_i)^\top$. If $\langle S,\mathcal PS\rangle>0$, then
\[
 \Pr\Bigl[\langle S,\mathcal PE\rangle<-\tfrac12\langle S,\mathcal PS\rangle\ \Bigm|\ \mathcal F\Bigr]\ \le\
 \exp\Bigl(-\frac{\langle S,\mathcal PS\rangle^2}{8\,\|\Sigma_E\|_{\mathrm{op}}\,\|\mathcal PS\|_F^2}\Bigr).
\]
\end{lemma}
\begin{proof}
By self-adjointness $\langle S,\mathcal PE\rangle=\langle\mathcal PS,E\rangle=\sum_ig_i\langle\mathcal PS,E_i\rangle$, which given $\mathcal F$ is centered normal with
variance $\operatorname{vec}(\mathcal PS)^\top\Sigma_E\operatorname{vec}(\mathcal PS)\le\|\Sigma_E\|_{\mathrm{op}}\|\mathcal PS\|_F^2$. Apply the Gaussian tail
$\Pr[Z<-u]\le e^{-u^2/(2\sigma^2)}$ with $u=\frac12\langle S,\mathcal PS\rangle$.
\end{proof}
\emph{Scope.} The lemma needs the interference to carry Gaussian randomness that is independent of $S$ and $\mathcal P$. In the Transformer,
$N$, $W_Q$ and $W_K$ are all functions of the same random weights, and they evolve along the flow, so this independence does not hold as stated.
Applying \cref{lem:cancel} to \cref{eq:interference} requires a decoupling step, for example conditioning on the task-relevant Gaussian
projections and a leave-one-out argument for the rest, and a stability estimate keeping the bound valid in time. Both are open.

\subsection{Proof obligations}\label{app:obligations}
\begin{center}\small
\begin{tabular}{@{}p{0.3\linewidth}p{0.3\linewidth}p{0.32\linewidth}@{}}
\toprule
Statement & Status & What would close it\\
\midrule
Seed in outcome regime (\cref{lem:seed}) & proved & ---\\
Outcome plateau from seed (\cref{prop:random-plateau}) & proved & ---\\
Alignment $\Rightarrow$ polynomial time and rollout (\cref{thm:process}) & proved & ---\\
Separability $\Rightarrow$ alignment (\cref{prop:separable}) & proved & ---\\
Matrix flow, score sensitivity (\cref{app:matrix}) & proved (exact identities) & ---\\
Regularity \ref{hyp:regular} & standard; ReLU switching needs a consistent flow & a.e.\ differentiability along the trajectory\\
Early alignment (bootstrap) & open & an interaction feature of inverse-polynomial size in polynomial time\\
Interference bound \cref{eq:interference} & open & decoupled version of \cref{lem:cancel}\\
Persistence of \ref{hyp:align} & open & time-uniform stability of features and interference\\
Mechanism recovery (routing, lookup) & not claimed & separate theorems on heads and MLP units\\
\bottomrule
\end{tabular}
\end{center}

\subsection{Diagnostics for the open hypotheses}\label{app:diagnostics}
The following measurements are planned; \emph{none has been run for this manuscript, and no result is reported}. They test whether
\ref{hyp:align} and \ref{hyp:sep} are plausible and how their constants scale; they cannot prove them. All use the executor architecture of
\cref{prop:represent} with every tensor trainable, momentum-free full-batch or small-step gradient descent as the approximation to flow
(adaptive optimizers reported separately), time normalized as the accumulated step size, and identical initial tensors for the process and
outcome runs.
\begin{enumerate}[leftmargin=*,itemsep=2pt]
\item \emph{Seeds.} Dense Gaussian; Gaussian plus a weak positional perturbation; plus a weak match-sensitive perturbation; plus both. Every
 seed must start with uniform predictions and nearly uniform relevant attention, so that it does not already count.
\item \emph{Alignment.} Along each run, record $a_{\mathrm{obs}}(t)=-\langle\nabla\Lproc,D(t)\rangle/(\Lproc\|D(t)\|)$ for the readout witness of
 \cref{prop:separable}, for the gradient of a rule-credit observable, and for a combined direction; report minima and lower quantiles over
 checkpoints, not means.
\item \emph{Separability.} Fit a minimum-norm witness $U$ on hidden states $h_\Phi(x)$ from one sample of histories and report
 $\gamma_{\mathrm{obs}}/\|U\|_F$ on held-out histories, including the four state--bit cases and all competing tokens.
\item \emph{Interference.} For each head, report $\langle S,\mathcal PE\rangle/\langle S,\mathcal PS\rangle$ from \cref{eq:interference} with a task subspace
 fixed in advance, at initialization and after substantial updates; the second tests the independence that \cref{lem:cancel} needs.
\item \emph{Mechanism.} Correct-source attention mass for each relevant head and the separation of the state--bit cases in the second MLP.
\item \emph{Scaling.} The first time at which free-running trace error falls below a threshold, for $n\in\{8,12,16,24,32,48,64\}$, with widths and
 scale set by the theorem's polynomial rule. Fits of $cn^a$ against $ce^{bn}$ are evidence for a conjecture, not a proof.
\end{enumerate}
